\documentclass[10.5pt]{article}
\usepackage[letterpaper, margin=0.6in, top=0.5in, bottom=0.5in]{geometry}
\usepackage{amsmath,amssymb,amsthm}
\usepackage{fancyhdr}
\usepackage{xcolor}
\usepackage{tikz}
\usepackage{enumitem}
\usepackage{tcolorbox}
\usepackage{hyperref}

% Define colors matching instructor board annotations
\definecolor{chalkblue}{RGB}{0, 105, 185}
\definecolor{chalkred}{RGB}{205, 35, 35}
\definecolor{chalkpink}{RGB}{230, 90, 120}
\definecolor{chalkcyan}{RGB}{15, 160, 190}
\definecolor{chalkgreen}{RGB}{20, 130, 50}
\definecolor{boxbg}{RGB}{252, 252, 254}
\definecolor{boxframe}{RGB}{150, 160, 180}

\pagestyle{fancy}
\fancyhf{}
\renewcommand{\headrulewidth}{0pt}
\renewcommand{\footrulewidth}{0pt}

\setlength{\parindent}{0pt}
\setlength{\parskip}{2.5pt}

\begin{document}

% ==============================================================================
% PAGE 1
% ==============================================================================
\noindent \textbf{\large Math 105L} \hfill \textbf{\large Lecture 5-2}\\[1pt]
{\textbf{\Large The Derivative Function \& Its Interpretation}}\\[-5pt]
\rule{\textwidth}{1.2pt}

\vspace{0.03in}

\textbf{\large 1. The Derivative as a Function}

Calculating the derivative point-by-point (e.g. at $a=1, 2, 4$) using limits involves repeating the exact same algebra each time. Instead, we calculate the derivative for an arbitrary variable $x$, creating the \textbf{derivative function} $f'(x)$:

\begin{tcolorbox}[colback=boxbg, colframe=boxframe, boxrule=0.8pt, arc=3pt, left=8pt, right=8pt, top=4pt, bottom=4pt]
\textbf{Definition of the Derivative Function:}\\[1pt]
Given a function $f(x)$, its derivative function is defined by the limit:
\[
  \mathbf{f'(x) = \lim_{h\to 0}\frac{f(x+h) - f(x)}{h}}
\]
\textbf{Equivalent Notations for the Derivative of $y = f(x)$:}
\[
  \mathbf{f'(x) = y' = \frac{dy}{dx} = \frac{d}{dx}\big[f(x)\big]}
\]
\end{tcolorbox}

\vspace{0.03in}

\textbf{\large 2. Qualitative Rules for Graphing $f'(x)$ from $f(x)$}

Since $f'(a)$ represents the slope of the tangent line to $y = f(x)$ at $x = a$:
\begin{itemize}[leftmargin=1.2em, topsep=1pt, itemsep=3pt]
  \item \textbf{$f(x)$ is increasing} $\implies$ slopes are positive $\implies \mathbf{f'(x) > 0}$ (graph of $f'$ lies \textbf{above} the $x$-axis).
  \item \textbf{$f(x)$ is decreasing} $\implies$ slopes are negative $\implies \mathbf{f'(x) < 0}$ (graph of $f'$ lies \textbf{below} the $x$-axis).
  \item \textbf{Local Extrema (Peaks and Valleys):} At turning points where $f(x)$ transitions from increasing to decreasing (or vice versa), the tangent line is horizontal $\implies$ slope is $0 \implies \mathbf{f'(a) = 0}$ (\textbf{$x$-intercepts} of $f'$).
  \item \textbf{Steepness:} Where $f(x)$ is steepest, $f'(x)$ reaches its maximum or minimum values. Where $f(x)$ flattens out, $f'(x)$ approaches $0$.
\end{itemize}

\vspace{0.04in}

\textbf{\large 3. Applied Interpretation \& Units of the Derivative}

\begin{tcolorbox}[colback=chalkblue!6, colframe=chalkblue, boxrule=1pt, arc=3pt, left=8pt, right=8pt, top=4pt, bottom=4pt]
\textbf{General Rule for Units:}\\[1pt]
The units of $f'(x)$ are always:
\[
  \mathbf{\text{Units of } f'(x) = \frac{\text{Units of output } (y)}{\text{Units of input } (x)}}
\]
\end{tcolorbox}

\vspace{0.03in}

\textbf{Worked Classroom Example: Eastbound Vehicle Motion}
\begin{itemize}[leftmargin=1.2em, topsep=1pt, itemsep=3pt]
  \item \textbf{Context:} Let $s(t)$ represent the distance due East from Durham of a car traveling along a straight highway, where $t$ is measured in hours and $s$ is measured in miles.
  \item \textbf{Units of $s'(t)$:}
  \[
    \text{Units} = \frac{\text{miles}}{\text{hour}} \quad \text{(instantaneous velocity)}
  \]
  \item \textbf{Physical Meaning of the Sign of $s'(t)$:}
  \begin{itemize}[label=--]
    \item $\mathbf{s'(t) > 0}$: Distance East is increasing $\implies$ The car is traveling \textbf{East} (away from Durham).
    \item $\mathbf{s'(t) < 0}$: Distance East is decreasing $\implies$ The car is traveling \textbf{West} (back toward Durham or westbound).
    \item $\mathbf{s'(t) = 0}$: The car is \textbf{instantaneously stationary} (changing direction or at rest).
  \end{itemize}
  \item \textbf{Synthesis:} Knowing the formula for $s'(t)$ gives the car's instantaneous speed and directional velocity at any exact hour $t$.
\end{itemize}

\end{document}

